\chapter{Introducción}

\section{Qué es el Sistema de Controlador RC con ESP32}

El sistema de Controlador RC con ESP32 te permite controlar proyectos electrónicos utilizando un teléfono Android y un microcontrolador ESP32.

La aplicación Android se comunica con el ESP32 mediante Bluetooth. Los comandos de control enviados desde el teléfono llegan al ESP32, que luego actualiza motores, servos, LEDs y otros dispositivos conectados a la placa.

Al mismo tiempo, el ESP32 puede enviar información de vuelta al teléfono, por ejemplo:

\begin{itemize}
    \item Valores de sensores
    \item Estados de botones
    \item Estado del sistema
    \item Telemetría en tiempo real
\end{itemize}

Esto crea un sistema de comunicación bidireccional entre la aplicación Android y el ESP32.

\section{Para Quién es Esta Guía}

Esta guía está dirigida a:

\begin{itemize}
    \item Estudiantes que están aprendiendo electrónica
    \item Aficionados a la robótica
    \item Makers que construyen proyectos DIY
    \item Principiantes que usan la plataforma ESP32
\end{itemize}

Solo se requiere conocimiento básico del IDE de Arduino y de cableado electrónico sencillo.

\section{Qué Puedes Construir}

Con el sistema de Controlador RC puedes construir muchos tipos de proyectos:

\begin{itemize}
    \item Autos RC
    \item Robots
    \item Sistemas de paneo para cámaras
    \item Plataformas de monitoreo de sensores
    \item Dispositivos controlados por Bluetooth
\end{itemize}

El firmware fue diseñado para ser modular, lo que permite a los usuarios ampliar el sistema con sensores adicionales, motores y dispositivos de control.